\section{Discussion}
\label{sec:discussion}

\subsection{Interpretation of Results}

Our five frameworks demonstrate that mathematical tools from topology, sheaf theory, information geometry, and spectral analysis provide genuinely new information about LLM evaluation---information inaccessible to standard metrics.

\para{Stability through topology.}
The topological drift detection framework provides something no existing evaluation metric offers: a formal guarantee that observed differences reflect genuine structural change rather than measurement noise.
The stability theorem (Theorem~\ref{thm:stability}) converts the Wasserstein distance of 1.77 into a provable lower bound on geometric displacement between model generations.
This is qualitatively different from reporting that ``accuracy increased by $X$\%''---it guarantees that the capability \emph{landscape} has changed, not just a single summary statistic.

\para{Composability through sheaves.}
The sheaf-theoretic framework formally characterizes a problem practitioners face intuitively: when can we trust aggregate scores?
Our consistency ratio directly measures information loss under aggregation.
\gemini 1.5 Pro's ratio of 1.83 means that reporting its category-level scores discards 83\% more information relative to signal than \llama 8B's uniform profile.
This provides actionable guidance: models with high consistency ratios require fine-grained evaluation; aggregate leaderboard rankings are formally misleading for these systems.

\para{Geometry beyond accuracy.}
The \fr distance reveals model relationships invisible to accuracy comparison.
The small distance (0.31) between \claude Opus and \gpt, despite different overall accuracies, indicates shared capability \emph{structure}---similar relative strengths and weaknesses across subjects.
This suggests architectural or training data similarities that raw performance numbers cannot capture.

\para{Phase transitions in failure.}
The emergence of $H_1$ features at threshold 0.8 reveals qualitative structure in the failure landscape.
Below this threshold, failures decompose hierarchically; above it, cyclic dependencies emerge.
This has practical implications: evaluation strategies that assume tree-structured failure modes (hierarchical decomposition, category-level analysis) break down for high-difficulty items.

\para{Spectral fingerprints.}
The spectral framework demonstrates that contamination creates detectable structural signatures in the performance matrix eigenspectrum.
Unlike surface-level methods that can be evaded by paraphrasing~\citep{yang2023rethinking}, spectral detection operates on the correlation structure of performance boosts, which is preserved even when individual items are modified.

\subsection{Limitations}

We identify several limitations that scope the applicability of our current results.

\para{Simulated data.}
Model performance profiles are constructed from published accuracy scores rather than measured directly via API calls.
While this provides realistic capability vectors, it does not capture the full distributional information (\eg per-question correctness patterns) that would strengthen the topological and spectral analyses.
Direct evaluation via model APIs is a necessary next step.

\para{Small sample size.}
With 12 models in $\sR^{57}$, the point cloud is sparse relative to ambient dimension.
The absence of $H_1$ features in the full cloud reflects this limitation rather than the absence of genuine topological structure.
Scaling to 100+ models using approximate persistence algorithms~\citep{sheehy2013linear} would enable richer topological analysis.

\para{Simplified sheaf construction.}
Our sheaf uses linear restriction maps (weighted averaging), which captures aggregation behavior but not the full complexity of evaluation composition.
Nonlinear restriction maps and computation of higher sheaf cohomology would provide stronger guarantees but require more sophisticated algebraic machinery.

\para{Contamination validation.}
We validate spectral detection using synthetic contamination injection rather than known real-world contamination cases.
Validation against documented contamination events~\citep{shi2023detecting, xu2024contamination} would strengthen the practical applicability of this framework.

\para{Computational cost.}
Persistent homology computation scales as $O(n^3)$ in the number of simplices, which grows exponentially with point cloud size.
For production-scale deployment with thousands of models, approximate algorithms or dimensionality reduction would be necessary.

\subsection{Broader Implications}

\para{Toward formal evaluation science.}
Our work suggests that LLM evaluation can and should be placed on formal mathematical foundations.
Just as statistical testing provides guarantees for hypothesis evaluation, topological and geometric tools can provide guarantees for model evaluation.
The stability theorem transforms evaluation from empirical measurement into provable characterization.

\para{Complementary to existing approaches.}
We emphasize that our frameworks \emph{complement} rather than replace existing evaluation methods.
Standard accuracy metrics remain essential for absolute performance assessment.
Our frameworks add a layer of structural analysis that detects when standard metrics are unreliable (high consistency ratios), when rankings are unstable (small persistence features), and when data integrity is compromised (spectral anomalies).

\para{Evaluation as mathematical object.}
By treating evaluation data as mathematical objects (point clouds, sheaves, statistical manifolds, random matrices), we open the door to importing decades of theoretical results from pure mathematics.
Each mathematical domain provides not just tools but \emph{theorems}---formal guarantees that the evaluation community currently lacks.
